\chapter{Pandemics}
\index{Pandemics}

\medskip
\noindent\textit{This chapter is a general overview. During an outbreak, follow current advice from public-health authorities and healthcare professionals in your area.}

\section*{What it is}
A pandemic is an epidemic that spreads across countries or continents and affects many people. The severity varies with the pathogen, immunity, available treatments, and the capacity of health systems. A pandemic may be caused by a new pathogen or by a known pathogen that spreads in a new way or reaches populations with limited immunity.

\section*{How pandemics are detected}
Public-health agencies use surveillance, laboratory testing, clinical reports, genomic analysis, and information from communities and healthcare services. Risk assessments change as evidence about transmission, severity, variants, treatments, and vaccines develops.

\section*{Reducing risk}
Measures depend on the pathogen and local recommendations. They may include vaccination when available, staying home when ill, ventilation, hand hygiene, masks or other protective equipment in specific settings, testing, treatment, contact management, and protection of people at higher risk. No single measure is effective in every situation.

\section*{During an outbreak}
Use reliable, current public-health information; check the date and source of claims; and avoid spreading unverified messages. Keep essential medicines available as advised, make a plan for caring for dependants, and seek help early if a chronic illness becomes difficult to manage. Public-health measures should be proportionate, evidence-informed, equitable, and respectful of rights.

\section*{When to Seek Medical Care}
Seek urgent care for severe breathing difficulty, chest pain, blue or grey lips, confusion, collapse, severe dehydration, or rapidly worsening illness. People at higher risk should ask a healthcare professional early about testing and treatment when they develop relevant symptoms.

\section*{Sources and further reading}
\begin{itemize}
\item World Health Organization. \textit{Preparing and preventing epidemics and pandemics}. \url{https://www.who.int/activities/preparing-and-preventing-epidemics-and-pandemics}
\item World Health Organization. \textit{Preparedness and Resilience for Emerging Threats (PRET)}. \url{https://www.who.int/initiatives/preparedness-and-resilience-for-emerging-threats}
\end{itemize}
